\documentclass[11pt]{amsart}

\usepackage[T1]{fontenc}
\usepackage{lmodern}
\usepackage{microtype}
\usepackage{mathtools,amssymb,amsthm}
\usepackage[hidelinks]{hyperref}

\hypersetup{
  pdftitle={Klein bottle sets of doubling 3|S|-3 refute a torsion-free
            small-doubling conjecture of Mohan and Neetu as stated}
}

\newtheorem{theorem}{Theorem}
\newtheorem{lemma}[theorem]{Lemma}
\newtheorem{corollary}[theorem]{Corollary}
\theoremstyle{remark}
\newtheorem*{remark}{Remark}
\newtheorem*{conjectureA}{Conjecture 6.3 (Mohan and Neetu)}

\title[Klein bottle sets of doubling $3|S|-3$]
{Klein Bottle Sets of Doubling $3|S|-3$ Refute a Torsion-Free
 Small-Doubling Conjecture of Mohan and Neetu as Stated}

\date{}

\subjclass[2020]{11P70, 20F60, 20F65, 11B13}
\keywords{small doubling, product set, torsion-free group, Klein bottle group,
Baumslag--Solitar group, right-ordered group, counterexample}

\begin{document}

\begin{abstract}
Let $S$ be a finite subset of a group $G$ and $S^{2}=\{s_{1}s_{2}:s_{1},s_{2}\in S\}$.
Conjecture~6.3 of Mohan and Neetu (arXiv:2607.11194v1) asserts that if $G$ is
torsion-free, $|S|\ge3$, $e\in S$ and $|S^{2}|\le3|S|-3$, then $\langle S\rangle$
is abelian. We exhibit counterexamples, at every admissible cardinality, inside
the Klein bottle group
$K=\mathbb{Z}\rtimes_{-1}\mathbb{Z}=\langle a,t\mid tat^{-1}=a^{-1}\rangle$:
writing $a^{m}t^{n}=(m,n)$ and
\[
  S_{k}=\{(0,0)\}\cup\{(i,1):1\le i\le k-1\}\qquad(k\ge3),
\]
one has $|S_{k}|=k$, $e\in S_{k}$, $|S_{k}^{2}|=3k-3=3|S_{k}|-3$ and
$\langle S_{k}\rangle=K$, which is nonabelian. Moreover $K$ carries a
right-invariant total order under which $e=\min S_{k}$. Freiman's torsion-free
conjecture at $|S^{2}|\le3|S|-4$ is not contradicted; the family sits one unit
above it.
\end{abstract}

\maketitle

\section{The conjecture}

Throughout, $G$ is a group with identity $e$, $S\subseteq G$ is finite, and
\[
  S^{2}=\{s_{1}s_{2}:s_{1},s_{2}\in S\}
\]
is the product set over \emph{ordered} pairs, so that $s^{2}\in S^{2}$ for every
$s\in S$ and $S\subseteq S^{2}$ as soon as $e\in S$. A set $S$ is called
\emph{abelian} when $\langle S\rangle$ is an abelian subgroup of $G$. In these
conventions, which are those of \cite{MN}, Freiman, Herzog, Longobardi and Maj
proved \cite[Theorem~1.3]{FHLM} that if $G$ is \emph{orderable}, $|S|\ge3$ and
$|S^{2}|\le3|S|-3$, then $S$ is abelian --- with no hypothesis $e\in S$. Mohan and
Neetu ask whether the hypothesis ``orderable'' can be weakened to
``torsion-free'' at the cost of adding $e\in S$. Verbatim
\cite[Conjecture~6.3]{MN}, quoted from the e-print arXiv:2607.11194v1 of 13 July
2026, its only version, where the statement stands in Section~6:

\begin{conjectureA}
Let $S$ be a nonempty finite subset of a torsion-free group $G$ with
$\left|S\right| \geq 3$ and $e\in S$. If $\left|S^{2}\right| \leq 3|S|-3$, then
$\langle S\rangle$ is an abelian subgroup of $G$.
\end{conjectureA}

\section{The counterexample}

Let $K$ be the Klein bottle group,
\[
  K=\mathbb{Z}\rtimes_{-1}\mathbb{Z}
   =\langle a,t\mid tat^{-1}=a^{-1}\rangle
   =\mathrm{BS}(1,-1),
\]
in the coordinates $a^{m}t^{n}\leftrightarrow(m,n)$, $m,n\in\mathbb{Z}$. The
multiplication and the inverse are
\begin{equation}\label{eq:law}
  (m,n)(p,q)=\bigl(m+(-1)^{n}p,\;n+q\bigr),
  \qquad
  (m,n)^{-1}=\bigl(-(-1)^{n}m,\,-n\bigr),
\end{equation}
with identity $e=(0,0)$.

\begin{lemma}\label{lem:tf}
$K$ is torsion-free, and it is not abelian.
\end{lemma}

\begin{proof}
By \eqref{eq:law} the second coordinate is additive, so $(m,n)^{r}$ has second
coordinate $rn$; if $n\ne0$ and $r\ne0$ this is nonzero and $(m,n)^{r}\ne e$. If
$n=0$ then $(m,0)^{r}=(rm,0)$, which is $\ne e$ when $m\ne0$ and $r\ne0$. Hence no
element other than $e$ has finite order. Finally $(1,0)(0,1)=(1,1)$ while
$(0,1)(1,0)=(-1,1)$, so $K$ is nonabelian.
\end{proof}

For an integer $k\ge3$ put $I=\{1,2,\dots,k-1\}$ and
\begin{equation}\label{eq:S}
  S_{k}=\{(0,0)\}\cup\{(i,1):i\in I\}
  =\{e,\;at,\;a^{2}t,\;\dots,\;a^{k-1}t\}.
\end{equation}

\begin{theorem}\label{thm:main}
For every integer $k\ge3$ the set $S_{k}\subseteq K$ of \eqref{eq:S} satisfies
\[
  |S_{k}|=k,\qquad e\in S_{k},\qquad
  S_{k}^{2}=S_{k}\cup\{(d,2):|d|\le k-2\},
\]
\[
  |S_{k}^{2}|=3k-3=3|S_{k}|-3,
\]
while $\langle S_{k}\rangle=K$ is nonabelian. Since $K$ is torsion-free by
Lemma~\ref{lem:tf}, Conjecture~6.3 of \cite{MN} is false as stated --- and false
for every admissible cardinality $|S|=k\ge3$, not at a single value.
\end{theorem}

\begin{proof}
The $k$ listed elements of \eqref{eq:S} are pairwise distinct because their
coordinates are, and $e=(0,0)\in S_{k}$.

Split $S_{k}^{2}$ by the second coordinate, which by \eqref{eq:law} adds and
therefore separates the three blocks below.
\begin{itemize}
\item Second coordinate $0$: the only product of two elements of $S_{k}$ with
      second coordinate $0$ is $e\cdot e=e$. One element.
\item Second coordinate $1$: the products are $e\cdot(i,1)=(i,1)=(i,1)\cdot e$ for
      $i\in I$. Exactly the $k-1$ elements $\{(i,1):i\in I\}$, and nothing new.
\item Second coordinate $2$: for $i,j\in I$,
      $(i,1)(j,1)=(i+(-1)^{1}j,2)=(i-j,2)$, so this block is
      $\{(d,2):d\in I-I\}$. As $I=\{1,\dots,k-1\}$ is an interval of integers,
      $I-I=\{-(k-2),\dots,k-2\}$ and $|I-I|=2k-3$ \emph{exactly}.
\end{itemize}
Hence $S_{k}^{2}=S_{k}\cup\{(d,2):|d|\le k-2\}$ and
\[
  |S_{k}^{2}|=1+(k-1)+(2k-3)=3k-3=3|S_{k}|-3,
\]
so the antecedent of Conjecture~6.3 holds, with equality.

For the consequent, $(1,1)(2,1)=(-1,2)$ and $(2,1)(1,1)=(1,2)$, so $S_{k}$ is not
even commuting elementwise. Moreover $\langle S_{k}\rangle=K$: by \eqref{eq:law}
$(1,1)^{-1}=(1,-1)$, and
\[
  (1,-1)(2,1)=(1-2,0)=(-1,0)=a^{-1},
  \quad
  (-1,0)(1,1)=(0,1)=t,
\]
so $\langle S_{k}\rangle$ contains $a$ and $t$, hence equals $K$, which is
nonabelian by Lemma~\ref{lem:tf}.
\end{proof}

\noindent
\textbf{The case $k=3$.} Here
$S_{3}=\{e,at,a^{2}t\}=\{(0,0),(1,1),(2,1)\}$, and the nine ordered products are
\[
\begin{array}{@{}l@{\quad}l@{\quad}l@{}}
 e\cdot e=(0,0), & e\cdot(1,1)=(1,1), & e\cdot(2,1)=(2,1),\\
 (1,1)\cdot e=(1,1), & (2,1)\cdot e=(2,1), & (1,1)(1,1)=(0,2),\\
 (2,1)(2,1)=(0,2), & (1,1)(2,1)=(-1,2), & (2,1)(1,1)=(1,2),
\end{array}
\]
with a single collision, $(at)^{2}=(a^{2}t)^{2}=t^{2}$. Thus
\[
  S_{3}^{2}=\{(0,0),(1,1),(2,1),(-1,2),(0,2),(1,2)\},
  \qquad
  |S_{3}^{2}|=6\le3\cdot3-3=6.
\]

\begin{corollary}\label{cor:ro}
Order $K$ by
\[
  (m,n)<(m',n')
  \quad\Longleftrightarrow\quad
  n<n',\ \text{or}\ n=n'\ \text{and}\ m<m'.
\]
This is a right-invariant total order on $K$, and $\min S_{k}=e$ for every $k\ge3$.
\end{corollary}

\begin{proof}
Right multiplication by $(p,q)$ sends $(m,n)$ to $(m+(-1)^{n}p,\,n+q)$. If $n<n'$
then $n+q<n'+q$. If $n=n'$ then both first coordinates are shifted by the same
integer $(-1)^{n}p$, so their order is preserved. Totality and transitivity are
those of the lexicographic order on $\mathbb{Z}^{2}$. Every element of
$S_{k}\setminus\{e\}$ has second coordinate $1>0$, so $e=\min S_{k}$.
\end{proof}

\begin{remark}[scope]
No minimality is claimed: nothing above asserts that $3|S|-3$ is the least
doubling attainable by an identity-containing finite set with nonabelian span in a
torsion-free group, nor that $k=3$ is the least cardinality in any other host
group. Freiman's torsion-free conjecture concerns $|S^{2}|\le3|S|-4$
\cite[\S1]{BPS}, and our family has $|S_{k}^{2}|=3k-3>3k-4$, one unit above that
threshold, so it is consistent with that conjecture. The host group is
$\mathrm{BS}(1,-1)$, which is the case excluded in \cite[Theorem~2.11]{MN}; what
is refuted is Conjecture~6.3 as printed, and nothing here bears on a version of it
in which that exclusion is in force. No priority or novelty is claimed for the
object: $K$ together with a three-element identity-containing set of nonabelian
span occurs in \cite[\S1]{BPS} and in \cite{AJ}, and the companion paper
\cite{MNS} is closed access and was not available to us.
\end{remark}

\section{Verification}

Every object the note claims is printed above: the group law \eqref{eq:law}, the
family \eqref{eq:S}, the nine products of the $k=3$ cell, the closed form for
$S_{k}^{2}$, the generating words for $a^{-1}$ and $t$, and the order of
Corollary~\ref{cor:ro}. A reader therefore needs no code. A short program
(\texttt{verify.py}, with its recorded output in \texttt{verify.output.txt})
accompanies the note anyway: it re-implements \eqref{eq:law} in exact integer
arithmetic, recomputes $S_{k}^{2}$ product by product for $3\le k\le30$ and
compares it with the closed form as a set of coordinate pairs rather than by
cardinality, re-derives the nine displayed products, and verifies
right-invariance of the order together with $\min S_{k}=e$; it also runs controls
of both polarities, so that neither of the two detectors it uses can be a
constant. All $304$ of its checks pass. It states in its own output what it does
not cover: torsion-freeness is confirmed structurally and over a bounded box only,
the general statement being Lemma~\ref{lem:tf}; no search for a smaller
counterexample is performed; and no bibliographic claim is checked.

\begin{thebibliography}{9}

\bibitem{AJ}
A.~Abdollahi and F.~Jafari,
\emph{Cardinality of product sets in torsion-free groups and applications in
group algebras},
J. Algebra Appl. \textbf{19} (2020), no.~5, 2050079;
\href{https://arxiv.org/abs/1808.08708}{arXiv:1808.08708}.

\bibitem{BPS}
K.~J. B\"or\"oczky, P.~P. P\'alfy, and O.~Serra,
\emph{On the cardinality of sumsets in torsion-free groups},
Bull. London Math. Soc. \textbf{44} (2012), no.~5, 1034--1041.
\href{https://doi.org/10.1112/blms/bds032}{doi:10.1112/blms/bds032};
\href{https://arxiv.org/abs/1009.6140}{arXiv:1009.6140}.

\bibitem{FHLM}
G.~A. Freiman, M.~Herzog, P.~Longobardi, and M.~Maj,
\emph{Small doubling in ordered groups},
J. Aust. Math. Soc. \textbf{96} (2014), no.~3, 316--325.
Theorem~1.3 is quoted as it is restated in \cite[Theorem~2.3]{MN}.

\bibitem{MN}
Mohan and Neetu,
\emph{On small doubling in right-ordered groups and Baumslag--Solitar groups~--~II},
arXiv:2607.11194v1, 2026.
\href{https://arxiv.org/abs/2607.11194}{arXiv:2607.11194}.
Statement numbers are quoted from this e-print.

\bibitem{MNS}
Mohan, Neetu, and Shankar,
\emph{On small doubling in right-ordered groups and Baumslag--Solitar groups},
Results Math. \textbf{81} (2026), no.~1, art.~12.
\href{https://doi.org/10.1007/s00025-025-02576-2}{doi:10.1007/s00025-025-02576-2}.
Closed access, no e-print.

\end{thebibliography}

\end{document}
